\section{Constructive source and effective outer-field inputs}\label{sec:effective-inputs}
Notes~44--45 \cite{n44,n45} make two parts of the fixed-gap contract
explicit: a finite positive source for any demanded absolute jet accuracy,
and complete-source outer field constants on a very large finite regime.
Neither alone instantiates the exterior collar and weighted-observable
constants in Theorem~\ref{thm:sj-transfer}.

\subsection{A finite model-gap template}
Fix $J=[95,97]$, $t_0=96$, $\ell=h=Y=1$ and $Y_b=4$. A candidate
center satisfies $\Theta_*(96,X)=n\pi$ on the established full-saddle
branch; no particular integer is selected. Write
\begin{align*}
 M(p)&=\frac{p(p-1)\pi^{-p/2}\Gamma(p/2)}{16},&
 \alpha&=M'/M,\\
 p&=w+\tfrac t2\alpha(p),&B&=1-\tfrac t2\alpha'(p),\\
 S_t(w)&=M(p)e^{-(p-w)^2/t}/\sqrt B.
\end{align*}
Here $B$ is the saddle determinant, distinct from the compatibility
derivative in Section~\ref{sec:spatial-join}. Suppose a full branch chart
supplies $b=\inf|B|>0$ and $a_j=\sup|\alpha^{(j)}(p)|$, $0\le j\le4$.
For $D=(\log S_t)_w$, direct implicit differentiation bounds
$D,D_w,D_{ww}$ by
\begin{align}
 D_0&=a_0+97a_2/(4b^2),\notag\\
 D_1&=a_1/b+97a_3/(4b^3)+97^2a_2^2/(4b^4),\notag\\
 D_2&=a_2/b^2+97a_1a_2/(2b^3)+97a_4/(4b^4)
                  +7\cdot97^2a_2a_3/(8b^5)+97^3a_2^3/(2b^6).
 \label{eq:fi-model-caps}
\end{align}
Consequently the carrier-relative jet constants in
\eqref{eq:sj-gap} are $1$, $D_0/(2L)$,
$(D_0^2+D_1)/(4L^2)$ and $(D_0^3+3D_0D_1+D_2)/(8L^3)$.
The time logarithmic derivative is exactly
\begin{equation}
 (\log S_t)_t=\frac{\alpha^2}{4}+\frac{\alpha'}{4B}
                            +\frac{t\alpha\alpha''}{8B^2},\qquad
 T_0=\frac{a_0^2}{4}+\frac{a_1}{4b}+\frac{97a_0a_2}{8b^2}.
 \label{eq:fi-time-cap}
\end{equation}
For compatible carrier logarithms
$\Theta=(\log S_+-\log S_-)/(2i)$, Taylor's integral formula gives
\begin{equation}
 |\Re\Theta(t,x+iy)-n\pi|\le\sigma D_0+\tau T_0+2D_1
 \quad(|x-X|\le2\sigma,\ |y|\le4).
 \label{eq:fi-phase}
\end{equation}
Indeed $\Theta$ is real on the real axis, its first transverse
displacement is imaginary, and $|\Theta_{zz}|\le D_1/4$.
The choices
\begin{equation}
 \sigma=\min\left(\tfrac12,\frac{\pi}{16(D_0+1)}\right),\qquad
 \tau\le\min\left(1,\frac{\pi}{16(T_0+1)}\right),\qquad D_1\le\pi/32
 \label{eq:fi-gap-choice}
\end{equation}
give phase deviation at most $3\pi/16<\pi/4$. Since
$|P|/A=|\cos\Theta|/\cosh(\Im\Theta)\ge|\cos(\Re\Theta)|$,
$c_0=1/2$ suffices. One must additionally supply collar injectivity,
separation and the root-speed gate \eqref{eq:sj-time}; a speed bound
is $V=T_0/\Omega_{\min}$ with
$\Omega_{\min}=\tfrac12\inf\Re D>0$ on real root charts.
The conservative time choice here can be smaller than $O(1/L)$.
These are finite sufficient inequalities, not an evaluated gap certificate.

The exact model residual gives
\begin{align}
 R_0&=97a_3/(16b^3)+5\cdot97^2a_2^2/(64b^4),\notag\\
 R_1&=97a_4/(16b^4)+97^2a_2a_3/(4b^5)
                              +5\cdot97^3a_2^3/(32b^6),\notag\\
 |E_P|&\le R_0A,\qquad |(E_P)_z|\le(D_0R_0+R_1)A/2.
 \label{eq:fi-residual-caps}
\end{align}
These do not infer field accuracy from residual size.

For a fully specified cutoff take
$\phi(s)=e^{-1/[s(1-s)]}$ on $(0,1)$ and zero elsewhere,
$b_*(s)=\int_{-\infty}^s\phi/\int\phi$. The integral over
$[1/4,3/4]$ is at least $e^{-16/3}/2$, and differentiation gives
\[
 \|b_*'\|_\infty\le2e^{4/3},\qquad
 \|b_*''\|_\infty\le32e^{4/3},\qquad
 \|b_*'''\|_\infty\le832e^{4/3}.
\]
Thus $\lambda=b_*((|x|-X+\sigma/2)/\sigma)$ is flat inside the
declared band, with these constants divided by $\sigma^j$.
Use $\eta_O=b_*((|x|-X-\sigma)/\sigma)$ for the heat partition.
The absolute value is harmless because both cutoffs are constant near
zero. Their exact normalizing integral is part of the analytic definition.

\subsection{Explicit compact demands and a clipped cubic source}
Let $R_c=X+2\sigma+3$, increased for any additional analytic buffer.
The real regulator saddle $L_s$ on $J$ obeys
$1/2<L_s<1$, $226<\kappa<618$, $0<\kappa'<13$ and
$1<M_0<e^{110}$. These follow by inserting $1/2,1$ in the real saddle
equation, then differentiating that equation. Also
$4e^{-4}<\Phi(0)\sqrt\pi/2<4$ from its first mode and complete
geometric majorant. Therefore a finite core floor is
\begin{equation}
 e_c=E_\eta e^{-R_c^2/452},\qquad
 E_\eta=\min((2/5)\eta e^{-4},1/618),\qquad
 \epsilon_g(t,x)\ge e_c\quad(t\in J,\ |x|\le R_c).
 \label{eq:fi-floor}
\end{equation}
The geometric interpolation defining $\epsilon_g$ is at least the
smaller endpoint lower bound. This floor is used only on the finite core.

For completeness, finite regulator caps consuming the full-support table
are
\begin{align}
 A_t&=4+(118+|\log\eta|)/8+R_c^2/512,&|r_t|&\le A_t,\notag\\
 R_1&\le R_c/113,&R_2&\le1/113,\notag\\
 T_1&\le R_c/256,&T_2&\le1/256,\notag\\
 T_h&\le R_c^2/14464,&c_r&\le1/113+R_c^2/113^2,\notag\\
 b_r&\le A_t+(1/113+R_c^2/113^2)/2.
 \label{eq:fi-regulator-caps}
\end{align}
Here $R_1,R_2$ again mean the regulator weights of
Section~\ref{sec:spatial-join}, not the model residual caps above.
To check these inequalities, use $|\chi'|<1/16$ on $J$,
$\beta\le1/226$, $|\beta'|<1/512$ and
$r=2\log E_0-\beta x^2$, retaining the switch amplitude mismatch.
Kernel norms through four can be bounded by
$k_\alpha\le a_{\alpha_x}a_{\alpha_y}$, with
\[
 (a_0,a_1,a_2,a_3,a_4)
   =(315/256,315/4,315,4725/4,16065/4).
\]
For source-moment caps take $U_0=1+\log111$ and, for $j=0,2$,
\[
 P_j=1+40\left\{
  \frac{U_0^{j+1}e^{97U_0^2+13U_0}}{j+1}
                         +\frac{\Gamma((j+1)/2)}2\right\}.
\]
Freeze $\bar d_0=\bar d_2=1$, $a_C=(2e_c)^{-1}$,
$b_C=e_c^{-2}$, $k_C=e_c^{-1}$ and declared common mean-gradient caps.
For $d_0,d_2\le\varepsilon\le1$ put $Z=\varepsilon/e_c^2$.
Then $s_C,w_C\le Z$,
$v_{m,C}\le[(1+m/2)P_2+1/2]Z$, and the pointwise time rows are at
most $(A_t+2P_2+1)Z$, $(A_t+3P_2+1)Z$ and
$(A_t+11P_2+3)Z$ respectively. Substitute these bounds and area at most
four into \eqref{eq:sj-curvature-ledger}--\eqref{eq:sj-source}.
If $D_r$ is the resulting coefficient after factoring $4Z$ from row
$r$, a sufficient common tolerance is
\begin{equation}
 \varepsilon\le\min\left(1,
        e_c^2\min_{D_r>0}\frac{\theta_{c,r}}{4D_r},
        \min_{0\le j\le3}\mu_sL^j A_{\min}\right).
 \label{eq:fi-source-demand}
\end{equation}
$D_r$ is an arithmetic expression in the displayed coefficients and
the fixed mean caps; exterior/gap debits remain outside this reserve.

\begin{proposition}[Positive clipped-cubic constructor]\label{prop:fi-source}
For every $0<\varepsilon\le1$ there is one nonnegative compact
piecewise-cubic source $f$, with rational coefficients and algebraic
splitting endpoints, whose exact heat field $G$ satisfies
\begin{equation}
 |\partial_z^j(G-H)|\le29\varepsilon/64
 \quad(0\le j\le3,\ t\in[95,97],\ x\in\R,\ |y|\le4).
 \label{eq:fi-source-error}
\end{equation}
An explicit construction is obtained by setting
\begin{align}
 Q_0&=\log(32/\varepsilon),&U&=\lceil4+\log(5+Q_0)\rceil,\notag\\
 F_0&=U^4e^{97U^2+4U},&B_4&=2^{20},\notag\\
 n_0&=\max\left(3,\left\lceil\sqrt{\log(2^{14}F_0/\varepsilon)/2}\right\rceil\right),
 &N&=n_0-1,\notag\\
 M&=\left\lceil U\max\{1,(B_4F_0/(48\varepsilon))^{1/4}\}\right\rceil,
 &h_s&=U/M,\notag\\
 b_s&=\lceil\log_2(64F_0/\varepsilon)\rceil.
 \label{eq:fi-recipe}
\end{align}
Use dyadic value and slope data for $\Phi_N$ at the rational nodes
$iU/M$, each accurate to $\varepsilon/(16F_0)$. Form the exact rational
Hermite cubic $p_i$ and set $f_i=\max(0,p_i)$; set $f=0$ above $U$.
Ceilings may use certified integer upper bounds, with a bounded extra
margin, to avoid deciding transcendental equality with an integer.
\end{proposition}
\begin{proof}
Since $\Phi\ge0$, clipping is a contraction:
\[
 |\max(0,p_i)-\Phi|\le|p_i-\Phi|
                 \le B_4h_s^4/384+a+h_sb/4,
\]
where $a,b$ bound nodal value/slope errors. A cubic has at most three
distinct real roots, so there are at most four polynomial or zero pieces
per cell. Rational coefficients give algebraic endpoints, including
multiple roots by exact polynomial algebra. The compact integral is
entire and exact heat despite corners in $f$, and uses the finite
Gaussian moments of Section~\ref{sec:hermite-source}.

The complete-source fourth derivative has the global bound
$\sup_{u\ge0}|\Phi^{(4)}(u)|<2^{20}$. To see the finite constant, use
the complete derivative-polynomial coefficients
\[
 (1875,30930,68096,41408,8448,512).
\]
The first-mode monomial envelope
is $75122271/80$; the geometric tail from $n=2$ is at most
$2055936470669686996992/28821512495674177$. Their sum is below $2^{20}$.
These are the rational majorants from the complete Hermite derivative
formula, not a finite-mode sample.

For $j\le3$, $\int_0^Uu^je^{97u^2+4u}\,du\le F_0$. The chosen
mesh makes interpolation cost at most $\varepsilon/8$. The complete
omitted nodal value/slope envelope is
\[
 \tau_0+h_s\tau_1/4\le2^{11}n_0^6e^{-3n_0^2}
        \le2^{11}e^{-2n_0^2}\le\varepsilon/(8F_0).
\]
The nodal arithmetic costs at most $5\varepsilon/64$ after integration.
For the complete tail, $U\ge4$ and the source-tail formula has
\[
 4\pi e^{4U}-194U-13-j/U\ge\pi e^{4U},\qquad
 97U^2+13U-\pi e^{4U}\le-2e^{4U}\le\log(\varepsilon/8),
\]
with prefactor at most one. It therefore costs at most
$\varepsilon/8$. Exact coefficient construction and clipping add zero.
The four paid categories sum to $29\varepsilon/64$.
\end{proof}
Nodal data are constructible by finite rational enclosures of $\pi$ and
exponentials; Note~44 gives an explicit, deliberately inefficient
Taylor/Machin prescription. No nodal evaluation or root isolation is
executed in this manuscript. Field evaluation, endpoint rounding and
continuous filtering are additional costs. With
$B_s=b_s+20+\lceil\log_2(M+1)\rceil$, a source-description upper bound
is $64(M+1)(B_s+1)$ bits. It includes algebraic root indices and isolating
intervals; it is not a running-time bound. As $q_0=\log(1/\varepsilon)$ grows,
\begin{equation}
 M=\varepsilon^{-1/4}\exp(O(\log^2q_0)),\qquad
 b_s,B_s=O(q_0+\log^2q_0).
 \label{eq:fi-cost}
\end{equation}
Thus a sufficient demand $\varepsilon=e^{-aX^2}$ can require
$M=\exp(aX^2/4+O(\log^2X))$ in this constructor. This is an upper
bound for this method, not a lower bound or an impossibility theorem.
Absolute jets uniform in $x$ do not by themselves give a global
shrinking-regulator observable comparison; the source is consumed on
the finite core and gap.

\subsection{A complete effective complex-buffer field estimate}
\begin{theorem}[Finite outer field and Cauchy-jet constants]\label{thm:effective-outer}
For the unchanged full saddle, $95\le t\le97$, $x=|\Re z|$ and
$L=\log(2+x)$,
\begin{equation}
 |H-P|\le12000A/x^2\quad(|\Im z|\le5,\ x\ge2^{2048}).
 \label{eq:ef-field}
\end{equation}
On $|\Im z|\le9/2$, $x\ge2^{2049}$,
\begin{equation}
 |\partial_z^j(H-P)|\le C_j A L^j/x^2\quad(0\le j\le3),\qquad
 (C_0,C_1,C_2,C_3)=(12000,384000,3072000,36864000).
 \label{eq:ef-jets}
\end{equation}
No value of $H$, gamma or a model gap is sampled to establish these
finite constants. The field theorem is independent of the regulator.
\end{theorem}
\begin{proof}
The first omitted digamma term with the complex remainder rule
\cite{DLMFGammaOuter2026} gives
\[
 \psi(v)=\Log v-\frac1{2v}+R(v),\qquad
 |R(v)|\le\frac{\sec^3(\arg(v)/2)}{12|v|^2}.
\]
For $\Re p\ge2$, $m=|p|\ge4096$, it follows that
\begin{equation}
 \alpha(p)=\tfrac12\Log(p/(2\pi))+r_0(p),\quad |r_0(p)|\le2/m,
 \qquad |\alpha^{(j)}(p)|\le j!/m^j\quad(1\le j\le4).
 \label{eq:ef-alpha}
\end{equation}
For the derivative bounds take a circle of radius $m/4$. It stays off
the negative axis and the rational poles, and there $|r_0|\le6/m$.
Cauchy gives $|r_0^{(j)}|\le6j!4^j/m^{j+1}$;
$1/(2j)+6\cdot4^j/4096<1$ pays the logarithmic term as well.

It suffices to work with $\lambda=\log x\ge1000$. For either arm,
after conjugation if needed, $w=\sigma_0+iT$ with
$-2\le\sigma_0\le3$, $T=x/2$. On the rectangle
\[
 \sigma_0+t\lambda/8\le\Re p\le\sigma_0+t\lambda/2,
 \qquad T\le\Im p\le T+t,
\]
one has $x/2\le|p|\le x$, and
\[
 (\lambda-3)/2-4/x\le\Re\alpha\le\lambda/2+4/x,
 \qquad 0<\Im\alpha<1.
\]
Thus the saddle map is an interior contraction with Lipschitz constant
$97/x<1/2$. Its unique holomorphic fixed point agrees with the established
full-saddle branch on overlapping charts. Writing $p-w=a+ih$ gives
\begin{equation}
 t\lambda/5\le a\le t\lambda/3,\quad 0<h<t/2,\quad
 |B-1|\le97/x,\quad 1/2<|B|<3/2,\quad\Re B>1/2.
 \label{eq:ef-saddle}
\end{equation}

Gaussian averaging of the canonical source, justified by its complete
superexponential tail, gives the exact convolution
\[
 H(t,z)=\frac1{\sqrt{\pi t}}\int_\R e^{-v^2/t}\xi(s+v)/8\,dv,
 \qquad s=(1+iz)/2.
\]
Split at real parts $2,-1$, reflect the left arm, and keep the length-three
middle interval. Each arm is
\begin{equation}
 I_w=\frac1{\sqrt{\pi t}}\int_{v_0}^\infty
              e^{-v^2/t}M(w+v)\zeta(w+v)\,dv,
 \qquad v_0=2-\sigma_0\in[-1,4].
 \label{eq:ef-arm}
\end{equation}
Shift to $v=r+ih$, retaining the oriented connector at $v_0$.
The real part of the gamma/zeta argument is at least two. Closing at
a finite right end and integrating the logarithmic derivative bound
shows its vertical face tends to zero against the Gaussian. The
following estimates also pay the actual infinite right tail.

For $v=p-w+u$ with $u\ge-a/2$, Taylor integral remainders give
\begin{align}
 \Gamma_p(u)&=\log M(p+u)-\log M(p)-\alpha(p)u,&
 |\Gamma_p(u)|&\le u^2/x,\notag\\
 E_p(u)&=\Gamma_p(u)-\alpha'(p)u^2/2,&
 |E_p(u)|&\le4|u|^3/(3x^2).
 \label{eq:ef-taylor}
\end{align}
The exact normalized integrand is $e^{-Bu^2/t+E_p(u)}$.
Both this and $e^{-Bu^2/t}$ have modulus at most $e^{-u^2/(2t)}$;
the line segment between their exponents has the same real-part bound.
The exponential difference integral and
$\int_\R|u|^3e^{-u^2/(2t)}du=4t^2$ therefore bound the central unit
arithmetic-mode error by $12t^{3/2}/x^2$, relative to $|S_t(w)|$.
Each missing Gaussian end and the actual unit-mode right end costs
at most $T_G=e^{-t\lambda^2/200}$; retain three such allowances.

The remaining contour costs are explicit. Put $A_r=\Re p=a+\sigma_0$
and $D_r=A_r/2+2\log(A_r/2)+2\le2A_r$ in this proof only.
Integrating \eqref{eq:ef-alpha} on the left horizontal segment gives
\[
 \int_{A_r-d}^{A_r}\Re\alpha(v+i(T+h))\,dv
                      \ge d\Re\alpha(p)-D_r.
\]
The bound $|\zeta|<2$ at real part at least two then gives the left
allowance below. The digamma series \cite{DLMFDigammaSeries2026} gives
$|\Im\alpha(2+iT)|<2$, paying the connector's factor $e^{2|h|}$.
For the middle, reflection, the pole-canceled Euler formula
\cite{DLMFZetaOuter2026} and the beta integral \cite{DLMFBeta2026} give
\begin{equation}
 \sup_{-1\le b\le2}|\xi(b+iT)/8|
                   \le225(1+|T|)|M(2+iT)|.
 \label{eq:ef-middle-bound}
\end{equation}
Indeed reflection reduces to $1/2\le b\le2$;
$|(q-1)\zeta(q)|\le|q|(1+2|q-1|)$ from Euler summation.
The beta integral bounds the gamma quotient by $\Gamma(b/2)<5$,
while $\pi^{1-b/2}<3$ and
$(2+|T|)^2/(1+T^2)\le5$. These give
$75(3+2|T|)\le225(1+|T|)$, including the canceled pole by continuity.
Finally the complete series has $|\zeta(q)-1|\le3\,2^{-\Re q}$.
Integration against the paid central Gaussian controls every non-unit
mode and its infinite right tail. The resulting per-arm relative ledger is
\begin{align}
 E_{\rm left}&=\frac{2|B|^{1/2}(a/2+1)}{\sqrt{\pi t}}
                                e^{-a^2/(4t)+D_r},\notag\\
 E_{\rm conn}&=\frac{2|B|^{1/2}|h|}{\sqrt{\pi t}}
                         e^{-(A_r-2)^2/t+D_r+2|h|},\notag\\
 E_{\rm mid}&=\frac{675|B|^{1/2}(1+T)}{\sqrt{\pi t}}
                         e^{16/t-(A_r-2)^2/t+D_r+2|h|},\notag\\
 E_{\rm arith}&=3\sqrt2|B|^{1/2}2^{-\sigma_0-a}e^{t(\log2)^2/2}.
 \label{eq:ef-paid-ledger}
\end{align}
The middle is counted once and can be compared to either carrier.
The complete zeta factor on the left and connector was already included;
no series estimate is extended outside its convergence half-plane.

For $\lambda\ge1000$, the finite reductions are
\begin{align}
 -a^2/(4t)+D_r&\le-95\lambda^2/100+(194/3)\lambda+6
                                           \le-\lambda^2/2,\notag\\
 T_G&\le x^{-3},& E_{\rm left}&\le66\lambda e^{-\lambda^2/2}\le x^{-3},\notag\\
 E_{\rm conn}&\le194e^{-\lambda^2/2+194}\le x^{-3},\notag\\
 E_{\rm mid}&\le1350x e^{-\lambda^2/2+195}\le x^{-3},\notag\\
 E_{\rm arith}&\le48e^{97/2-95\lambda/10}\le x^{-3}.
 \label{eq:ef-finite-decay}
\end{align}
They follow from \eqref{eq:ef-saddle}, $A_r-2\ge a/2$ and elementary
exponential inequalities; the quadratic margins are positive at $1000$
and increasing thereafter. Since $t^{3/2}<970$, both arms, all three
Gaussian/unit ends per arm, both connectors, all arithmetic modes and
the one middle give
\[
 |H-P|/A\le11640/x^2+7/x^3<12000/x^2.
\]
The finite threshold $2^{2048}$ lies within this regime because
$\log2>1/2$.

For jets use the exact identity
$D=\alpha+t\alpha''/(4B^2)$, which gives $|D|\le\log x$.
At $|\Im z|\le9/2$, $x\ge2^{2049}$, the full circle of radius
$\rho=1/(4L)$ stays in $|\Im z|\le5$, with height at least $x-1$.
Each carrier has $z$ logarithmic derivative at most $L$, so its modulus
on that circle is less than twice its central modulus. Cauchy's formula
\cite{DLMFCauchy1931} yields
\[
 \frac{|\partial_z^j(H-P)|}{A}
   \le\frac{24000j!(4L)^j}{(x-1)^2}
   \le\frac{96000j!4^jL^j}{x^2}.
\]
Retain the sharper $C_0$ from the field estimate. This proves the stated
table without differentiating an asymptotic remainder.
\end{proof}

On a gap chart with $X\ge2^{2050}$,
$||\Re z|-X|\le2$, $|\Im z|\le9/2$, one has local height at least
$X/2$ and local logarithm at most $2\log(2+X)$. Hence the formerly
uninstantiated field input in \eqref{eq:sj-source-gap} is paid by
\begin{equation}
 \mu_o\le\frac{32C_3}{X^2}
             =\frac{1179648000}{X^2}\le\frac1{16}.
 \label{eq:ef-gap-consumer}
\end{equation}
With $\mu_s\le1/16$, the $c_0=1/2$ floor has strict margin.
This is a uniform full-buffer theorem on $[95,97]$, not an accepted
phase index, finite source, full observable budget or practical onset.
Section~\ref{sec:uniform-assembly} supplies the finite collar/unit,
signed-pole and weighted-exterior constants and consumes them in a
complete budget. It also extends this field proof to $[63,97]$ by
rechecking the changed lower-time contour inequalities. The separate all-height late-time table is proved in
Section~\ref{sec:effective-positive}; the subsequent compact extension
and positive attachment are supplied in Section~\ref{sec:all-time}.
